\documentclass[11pt]{article}
\usepackage[T1]{fontenc}
\usepackage{lmodern,amsmath,amssymb,amsthm,mathtools}
\usepackage[margin=1in]{geometry}
\usepackage[colorlinks=true,linkcolor=blue,citecolor=blue,urlcolor=blue]{hyperref}
\newtheorem{theorem}{Theorem}
\newtheorem{lemma}[theorem]{Lemma}
\newtheorem{corollary}[theorem]{Corollary}
\theoremstyle{remark}
\newtheorem{remark}[theorem]{Remark}
\DeclareMathOperator{\diag}{diag}
\DeclareMathOperator{\conv}{conv}
\DeclareMathOperator{\lcm}{lcm}
\DeclareMathOperator{\Tor}{Tor}
\DeclareMathOperator{\adj}{adj}
\newcommand{\Z}{\mathbb Z}
\newcommand{\R}{\mathbb R}
\newcommand{\K}{\mathcal K}
\newcommand{\one}{\mathbf 1}
\begin{document}
\begin{center}
{\Large\bfseries Sharp orders of arithmetical critical groups}\par\vspace{5pt}
{\large via canonical lattice simplices}\par\vspace{8pt}
September 2026
\end{center}
\begin{abstract}
We identify the critical group of an arithmetical structure on a connected simple graph with the vertex-lattice quotient of a canonical lattice simplex. Applying the sharp multiplicity theorem of Averkov, Kasprzyk, Lehmann and Nill, we determine the largest possible critical-group order among all such structures on graphs with a fixed number of vertices. Complete graphs attain the bound. The same bound holds for stars with a fixed number of leaves, proving the largest-order conjecture of Archer, Diaz-Lopez, Glass and Louwsma. We also determine the maximizing arithmetical structures on complete graphs and stars.
\end{abstract}

\section{Introduction and main result}

Let $G$ be a connected simple graph on $[n]=\{1,\ldots,n\}$, with $n\geq2$, and let $A(G)$ be its adjacency matrix. An \emph{arithmetical structure} on $G$ is a pair of positive integer vectors $(d,r)$ such that $r$ is primitive and
\[
 Lr=0,\qquad L=\diag(d)-A(G).
\]
Its critical group is the finite abelian group
\[
 \K(G;d,r)=\Tor(\Z^n/L\Z^n).
\]
The ordinary graph Laplacian is the case $r=\one$. Arithmetical structures allow other diagonal entries, producing critical groups whose orders can be much larger than the numbers of spanning trees of the underlying graphs.

Write $s_1=2$ and $s_{j+1}=s_1\cdots s_j+1$ for Sylvester's sequence, and define
\begin{equation}\label{eq:bound}
 M_n=\begin{cases}
 n^{n-2},&2\leq n\leq4,\\
 128,&n=5,\\
 3(s_{n-2}-1)^2,&n\geq6.
 \end{cases}
\end{equation}
For stars we use the convention that $S_n$ has $n$ leaves and hence $n+1$ vertices.

\begin{theorem}\label{thm:main}
For every arithmetical structure on a connected simple graph $G$ with $n\geq2$ vertices,
\[
 |\K(G;d,r)|\leq M_n.
\]
For each $n$, the bound is attained by an arithmetical structure on the complete graph $K_n$. Moreover, $M_n$ is the largest critical-group order among arithmetical structures on the star $S_n$.
\end{theorem}

Archer, Diaz-Lopez, Glass and Louwsma~\cite{ADGL} determine critical groups of arithmetical structures on stars and complete graphs. Their Conjecture~24 asserts the last conclusion of Theorem~\ref{thm:main} for $n\geq6$, together with its complete-graph counterpart. Their notation is $a_1=1$ and $a_{j+1}=a_j^2+a_j$, so that $a_j=s_j-1$. Thus Theorem~\ref{thm:main} proves that conjecture and gives a sharp bound over all connected simple graphs with a fixed number of vertices.

The geometric input is the sharp canonical-simplex multiplicity bound of Averkov, Kasprzyk, Lehmann and Nill~\cite{AKLN}. We use their theorem, including its equality classification, as an external result; we do not give a new proof of that bound. Taking the convex hull of Laplacian columns is the earlier construction of Braun and Meyer~\cite{BM}. Balletti, Hibi, Meyer and Tsuchiya~\cite{BHMT} show that ordinary Laplacian simplices of simple strongly connected digraphs have a unique interior lattice point. Our observation is that a symmetric arithmetical Laplacian gives such a simplex in the saturated lattice perpendicular to its positive kernel vector, and that its critical group is exactly the vertex-lattice quotient. This makes the multiplicity theorem applicable to the graph-order problem.

\section{The simplex and its vertex-lattice quotient}

For a rank-$m$ lattice $N$, a full-dimensional lattice simplex $P\subset N\otimes_\Z\R$ is called \emph{canonical} if its unique interior lattice point is the origin. Its \emph{multiplicity} is the index
\[
 \operatorname{mult}_N(P)
 =\left[N:\sum_{v\text{ a vertex of }P}\Z v\right].
\]
In particular, the ambient lattice is part of this definition.

\begin{lemma}\label{lem:bridge}
Let $(d,r)$ be an arithmetical structure on a connected simple graph with $n\geq2$ vertices. Let $v_1,\ldots,v_n$ be the columns of $L=\diag(d)-A(G)$ and set
\[
 N=\{x\in\Z^n:r^{\mathsf T}x=0\},\qquad
 P=\conv(v_1,\ldots,v_n).
\]
Then $P$ is a canonical $(n-1)$-dimensional simplex in $N$, and
\[
 N/\sum_{i=1}^n\Z v_i\cong\K(G;d,r).
\]
\end{lemma}
\begin{proof}
Put $R=\diag(r)$. The matrix $RLR$ is the weighted Laplacian with edge weights $r_i r_j$, because
\[
 r_i^2d_i=\sum_{j:\{i,j\}\in E(G)}r_i r_j.
\]
Consequently, for every $z\in\R^n$,
\[
 z^{\mathsf T}RLRz
 =\sum_{\{i,j\}\in E(G)}r_i r_j(z_i-z_j)^2.
\]
Connectedness implies that $L$ has rank $n-1$ and $\ker_{\R}L=\R r$. Symmetry puts each $v_i$ in $N$. The only linear relation among the columns is a multiple of $\sum_i r_i v_i=0$. Since $\sum_i r_i>0$, no nonzero relation is an affine relation, so the columns are affinely independent. The positive coefficients $r_i/\sum_jr_j$ express $0$ as an interior point of their simplex.

Every coordinate of every $v_i$ is at least $-1$. For each $j$, column $v_j$ has positive $j$th coordinate $d_j$: every vertex has a neighbor, and $d_jr_j$ is the positive sum of the neighboring $r$-entries. An interior point of $P$ uses all vertices with strictly positive coefficients. Each of its coordinates is therefore strictly greater than $-1$. An interior lattice point $x$ has all $x_j\geq0$; the equation $r^{\mathsf T}x=0$ and positivity of $r$ force $x=0$. This proves that $P$ is canonical. Its vertices are primitive, since each column has an entry $-1$ at a neighboring vertex.

Finally, primitivity of $r$ gives an exact sequence
\[
 0\longrightarrow N/L\Z^n
 \longrightarrow\Z^n/L\Z^n
 \xrightarrow{\;r^{\mathsf T}\;}\Z
 \longrightarrow0.
\]
The left term is finite because $\operatorname{rank}L=n-1$. It equals the torsion subgroup of the middle term: torsion maps to zero in $\Z$, and every element of the finite kernel is torsion. The columns generate $L\Z^n$, giving the claimed isomorphism.
\end{proof}

We state the precise geometric input in lattice language.

\begin{theorem}[Averkov--Kasprzyk--Lehmann--Nill]\label{thm:AKLN}
Let $P$ be a canonical $m$-dimensional lattice simplex, where $m\geq1$. Then
\[
 \operatorname{mult}(P)\leq M_{m+1}.
\]
Up to a lattice isomorphism, the simplex attaining equality is unique. It is
\[
 \begin{cases}
 \conv(0,(m+1)e_1,\ldots,(m+1)e_m)-\one,&1\leq m\leq3,\\[2pt]
 \conv(0,2e_1,8e_2,8e_3,8e_4)-\one,&m=4,\\[2pt]
 \conv(0,s_1e_1,\ldots,s_{m-2}e_{m-2},3te_{m-1},3te_m)-\one,
 &m\geq5,
 \end{cases}
\]
where $t=s_{m-1}-1$ in the last line and the displayed models have ambient lattice $\Z^m$.
\end{theorem}

This is Proposition~2.8 and Theorem~2.9 of~\cite{AKLN}. Combining it with Lemma~\ref{lem:bridge}, with $m=n-1$, proves the upper bound in Theorem~\ref{thm:main}.

\section{Attainment and the star bound}

We record an elementary order formula. If $L$ is a symmetric positive semidefinite integral matrix of rank $n-1$ with primitive null vector $r$, then
\begin{equation}\label{eq:adjugate}
 \adj(L)=c\,rr^{\mathsf T},\qquad
 c=|\Tor(\Z^n/L\Z^n)|.
\end{equation}
Indeed the rank-one adjugate is a rational multiple of $rr^{\mathsf T}$. All its entries are integral, and $\gcd_{i,j}(r_i r_j)=1$, so the multiple is an integer. The gcd of the $(n-1)$-minors is the order of the torsion cokernel by Smith normal form. Positivity of the nonzero eigenvalues fixes the positive sign in~\eqref{eq:adjugate}.

Let $q_1,\ldots,q_n$ be positive integers with
\begin{equation}\label{eq:unit}
 \sum_{i=1}^n\frac1{q_i}=1.
\end{equation}
Put $\ell=\lcm(q_1,\ldots,q_n)$, $r_i=\ell/q_i$ and $d_i=q_i-1$. All $q_i\geq2$, and $r$ is primitive: for each prime dividing $\ell$, some denominator has its full valuation. Since $\sum_i r_i=\ell$, these data define an arithmetical structure on $K_n$ with
\[
 L=\diag(q_1,\ldots,q_n)-J,
\]
where $J$ is the all-ones matrix. Deleting row and column $j$ and applying the determinant lemma gives the principal cofactor
\[
 \prod_{i\ne j}q_i\left(1-\sum_{i\ne j}\frac1{q_i}\right)
 =\frac{\prod_i q_i}{q_j^2}.
\]
By~\eqref{eq:adjugate},
\begin{equation}\label{eq:order}
 |\K(K_n;d,r)|=\frac{\prod_iq_i}{\ell^2}.
\end{equation}
This is also the order consequence of the group formula in~\cite{ADGL}.

For $2\leq n\leq4$, take $q_i=n$ for every $i$; formula~\eqref{eq:order} gives $M_n=n^{n-2}$. For $n=5$, take
\[
 q=(2,8,8,8,8),
\]
giving order $128$. For $n\geq6$, put
\begin{equation}\label{eq:extremal}
 t=s_{n-2}-1=\prod_{i=1}^{n-3}s_i,\qquad
 q=(s_1,\ldots,s_{n-3},3t,3t,3t).
\end{equation}
The telescoping identity $\sum_{i=1}^{n-3}1/s_i=1-1/t$ proves~\eqref{eq:unit}. Here $\ell=3t$, and~\eqref{eq:order} gives
\[
 \frac{t(3t)^3}{(3t)^2}=3t^2=M_n.
\]
These are the extremal constructions already proposed in~\cite{ADGL}; their optimality follows from the upper bound.

Now consider an arithmetical structure on $S_n$. Denote its center by $0$, put $q_i=d_i$ at its leaves, and let $d_0$ be its central label. The arithmetical equations give $r_i=r_0/q_i$ and
\[
 \sum_{i=1}^n\frac1{q_i}=d_0.
\]
Primitivity implies $r_0=\ell=\lcm(q_1,\ldots,q_n)$. The principal cofactor obtained by deleting the center is $\prod_iq_i$, so~\eqref{eq:adjugate} gives
\begin{equation}\label{eq:starorder}
 |\K(S_n;d,r)|=\frac{\prod_iq_i}{\ell^2}.
\end{equation}
Replace each $q_i$ by $q'_i=d_0q_i$. The new tuple satisfies~\eqref{eq:unit}, has least common multiple $d_0\ell$, and determines a complete-graph structure of order
\[
 \frac{\prod_i q'_i}{(d_0\ell)^2}
 =d_0^{n-2}|\K(S_n;d,r)|.
\]
The complete-graph bound therefore proves the stronger inequality
\begin{equation}\label{eq:center}
 d_0^{n-2}|\K(S_n;d,r)|\leq M_n.
\end{equation}
This scaling is the order calculation underlying~\cite[Proposition~20]{ADGL}. Conversely, each maximizing unit-fraction tuple above defines a star structure with $d_0=1$ and order $M_n$ by~\eqref{eq:starorder}. This completes the proof of Theorem~\ref{thm:main}.

\section{Maximizing structures on complete graphs and stars}

The equality classification in the geometric theorem gives more than the numerical answer to the conjecture.

\begin{corollary}\label{cor:equality}
Up to permutation of vertices, the maximizing arithmetical structure on $K_n$ is unique for each $n\geq2$. It is given by $d_i=q_i-1$ and $r_i=\ell/q_i$, where $\ell=\lcm(q_i)$ and
\[
 q=\begin{cases}
 (n,\ldots,n),&2\leq n\leq4,\\
 (2,8,8,8,8),&n=5,\\
 (s_1,\ldots,s_{n-3},3t,3t,3t),&n\geq6,
 \end{cases}
\]
with $t=s_{n-2}-1$. For $S_n$ with $n\geq3$, the maximizing structure is likewise unique up to permutation of leaves: its central label is $d_0=1$, its leaf labels are the displayed $q_i$, and $(r_0,r_1,\ldots,r_n)=(\ell,\ell/q_1,\ldots,\ell/q_n)$.
\end{corollary}
\begin{proof}
For a complete-graph structure put $q_i=d_i+1$ and $h=\sum_i r_i$. The equations give $q_i r_i=h$, hence $\sum_i1/q_i=1$. The barycentric coordinates of $0$ in the simplex of Lemma~\ref{lem:bridge} are $r_i/h=1/q_i$. Lattice isomorphisms preserve these coordinates up to permutation. The equality models in Theorem~\ref{thm:AKLN} have reciprocal barycentric coordinates precisely the tuples in the statement. For the last model, the nonzero-axis vertices have coordinates $1/s_1,\ldots,1/s_{n-3},1/(3t),1/(3t)$, and the remaining coordinate is $1/(3t)$ by the telescoping identity. The low-dimensional models give the other two tuples. Thus equality determines $q$ up to permutation. It then determines $d$ and the primitive vector $r$ uniquely.

For a star with $n\geq3$, equality in~\eqref{eq:center} forces $d_0=1$. Its leaf tuple therefore determines a complete-graph structure of the same order, so the first part applies. The displayed constructions establish existence in every case.
\end{proof}

For $S_2$ every arithmetical critical group is trivial, and uniqueness is not asserted. Nor do we classify all graphs attaining the universal bound in Theorem~\ref{thm:main}. The numerical bound and its attainment need no such classification.

\begin{thebibliography}{9}

\bibitem{ADGL}
K.~Archer, A.~Diaz-Lopez, D.~Glass, and J.~Louwsma,
\emph{Critical groups of arithmetical structures on star graphs and complete graphs},
Electron. J. Combin. \textbf{31} (2024), no.~1, Paper P1.5.
\href{https://doi.org/10.37236/11729}{doi:10.37236/11729}.

\bibitem{AKLN}
G.~Averkov, A.~Kasprzyk, M.~Lehmann, and B.~Nill,
\emph{Sharp bounds on fake weighted projective spaces with canonical singularities},
in \emph{Varieties, Polyhedra, Computation}, EMS Series of Congress Reports,
EMS Press, 2025, pp.~319--334.
\href{https://doi.org/10.4171/ecr/22/10}{doi:10.4171/ecr/22/10}.
Preprint: \href{https://arxiv.org/abs/2105.09635}{arXiv:2105.09635}.

\bibitem{BHMT}
G.~Balletti, T.~Hibi, M.~Meyer, and A.~Tsuchiya,
\emph{Laplacian simplices associated to digraphs},
Ark. Mat. \textbf{56} (2018), 243--264.
\href{https://doi.org/10.4310/ARKIV.2018.v56.n2.a3}{doi:10.4310/ARKIV.2018.v56.n2.a3}.

\bibitem{BM}
B.~Braun and M.~Meyer,
\emph{Laplacian simplices},
Adv. Appl. Math. \textbf{114} (2020), 101976.
\href{https://doi.org/10.1016/j.aam.2019.101976}{doi:10.1016/j.aam.2019.101976}.

\end{thebibliography}
\end{document}
